\documentclass[11pt]{article}
\usepackage[margin=1in]{geometry}
\usepackage{amsmath,amssymb,amsthm}
\usepackage{graphicx}
\usepackage{hyperref}

\newcommand{\M}{\mathcal M}
\newcommand{\h}{\mathsf h}
\newcommand{\bmo}{\mathsf{bmo}}

\newtheorem{theorem}{Theorem}
\newtheorem{lemma}{Lemma}
\newtheorem{assumption}{Conditional Dependency}

\title{A conditional reduction of the endpoint complex\\
Hardy--BMO problem to the \(h_\infty\) and conditioned endpoints}
\author{Automatic Math Discovery Run \texttt{fa\_banach\_001}}
\date{June 16, 2026}

\begin{document}
\maketitle

\section*{Status}

This is a conditional reduction for the open problem in
Randrianantoanina's paper \cite{Ran2021}. It does not claim a full solution.
It proves that, in the range \(r>1\), the missing reverse inclusion for the
endpoint \(bmo^c\) complex interpolation problem follows from the corresponding
reverse inclusion for the endpoint \(h_\infty^c\) complex interpolation
problem. It also records a further sufficient reduction to an endpoint
interpolation statement for Junge's conditioned sequence spaces. The
\(h_\infty^c\) complex endpoint is explicitly left open in the later paper
\cite{Ran2022}.

\section*{Source problem}

The source paper proves that, for finite indices \(0<p,q<\infty\),
\[
[\h_p^c(\M),\h_q^c(\M)]_\theta=\h_r^c(\M),
\qquad
  \frac1r=\frac{1-\theta}{p}+\frac{\theta}{q}.
\]
Immediately after that theorem, it asks whether the corresponding statement
remains valid when \(0<p<1\) and the endpoint couple
\((\h_p^c(\M),\bmo^c(\M))\) is used.

\begin{center}
\includegraphics[width=\textwidth]{figures/open_problem_crop.png}
\end{center}

\noindent
Thus the desired full endpoint identity is
\[
[\h_p^c(\M),\bmo^c(\M)]_\theta=\h_r^c(\M),
\qquad
\frac1r=\frac{1-\theta}{p}.
\]

\section*{Conditional dependencies}

The first unproved dependency in the reduction is the following endpoint
reverse inclusion.

\begin{assumption}[\(h_\infty^c\) complex endpoint reverse inclusion]
Let \(0<p<1\), \(0<\theta<1\), and \(1/r=(1-\theta)/p\), with \(r>1\). Assume
that
\[
\h_r^c(\M)\hookrightarrow [\h_p^c(\M),\h_\infty^c(\M)]_\theta
\]
continuously.
\end{assumption}

The full complex identity
\[
[\h_p^c(\M),\h_\infty^c(\M)]_\theta=\h_r^c(\M)
\]
would imply this dependency. Randrianantoanina \cite{Ran2022} proves the
corresponding real/K-method theorem for \((\h_p^c,\h_\infty^c)\), but notes
that the complex version remains open.

This packet also isolates a lower-level sufficient dependency in the
conditioned sequence spaces. Let \(L_s^{\rm cond}(\M;\ell_2^c)\) denote
Junge's column conditioned space, with the \(L_\infty\)-endpoint framework as
in \cite{Ran2022}. Let
\[
\Pi(a)=\sum_{n\ge1}\big(\mathcal E_n(a_n)-\mathcal E_{n-1}(a_n)\big)
\]
be the natural martingale-difference projection, initially defined on finite
sequences. By the Kadison--Schwarz calculation in \cite[Lemma 2.2]{Ran2021},
\(\Pi\) is contractive on finite conditioned square functions.

\begin{assumption}[conditioned \(L_\infty\) endpoint interpolation]
Let \(0<p<1\), \(0<\theta<1\), and \(1/r=(1-\theta)/p\), with \(r>1\). Assume
that the endpoint conditioned couple is compatible with \(\Pi\) and satisfies
\[
L_r^{\rm cond}(\M;\ell_2^c)\hookrightarrow
[L_p^{\rm cond}(\M;\ell_2^c),
 L_\infty^{\rm cond}(\M;\ell_2^c)]_\theta
\]
continuously. Here compatibility means that the finite-sequence projection
\(\Pi\) extends boundedly at the \(L_\infty^{\rm cond}\) endpoint onto
\(\h_\infty^c(\M)\), as it does for \(0<s<\infty\) in
\cite[Lemma 2.2]{Ran2021}.
\end{assumption}

\section*{Idea of the proof}

The previous one-sided packet used Wolff's inclusion lemma to prove
\[
[\h_p^c(\M),\bmo^c(\M)]_\theta\hookrightarrow \h_r^c(\M)
\]
for \(r>1\). To get equality, only the reverse inclusion is missing. The
stronger endpoint \(h_\infty^c\) sits continuously inside \(bmo^c\): bounded
full conditioned square function controls the conditioned BMO tail norms.
Complex interpolation is monotone under continuous embeddings of endpoint
couples. Therefore a reverse inclusion through \(h_\infty^c\) immediately
passes to the larger endpoint \(bmo^c\).

The additional conditioned-space reduction uses the same projection \(\Pi\).
If the martingale difference sequence \(D_cx=(dx_n)\) of an element
\(x\in\h_r^c\) has an admissible analytic representation in the conditioned
couple, applying \(\Pi\) pointwise to that analytic family produces an
admissible analytic representation in the Hardy couple and keeps the value
\(\Pi D_cx=x\).

\begin{lemma}
There is a continuous embedding
\[
\h_\infty^c(\M)\hookrightarrow \bmo^c(\M).
\]
\end{lemma}

\begin{proof}
For a finite martingale \(x\), the \(bmo^c\) tail seminorm is controlled by
the bounded conditioned square function. Indeed, for \(m>n\), martingale
orthogonality gives
\[
\mathcal E_n\left(|x_m-x_n|^2\right)
  =\mathcal E_n\left(\sum_{k=n+1}^m |dx_k|^2\right)
  =\mathcal E_n\left(\sum_{k=n+1}^m
       \mathcal E_{k-1}(|dx_k|^2)\right).
\]
The last positive operator is dominated in norm by
\(\|s_c(x)^2\|_\infty\). The first-coordinate term is also controlled by the
standard \(h_\infty^c\) norm convention. Passing from finite martingales to
the completed endpoint gives the continuous embedding.
\end{proof}

\begin{theorem}
Let \(0<p<1\), \(0<\theta<1\), and put
\[
\frac1r=\frac{1-\theta}{p}.
\]
Assume \(r>1\) and assume the \(h_\infty^c\) complex endpoint reverse
inclusion above. Then
\[
[\h_p^c(\M),\bmo^c(\M)]_\theta=\h_r^c(\M)
\]
with equivalent quasi-norms.
\end{theorem}

\begin{proof}
The previous one-sided result proves
\[
[\h_p^c(\M),\bmo^c(\M)]_\theta\hookrightarrow \h_r^c(\M)
\]
when \(r>1\). For completeness, we recall the mechanism. Choose
\(\eta\in(0,1)\) with
\[
1=\frac{1-\eta}{p}+\frac{\eta}{r}.
\]
The finite-index complex interpolation theorem of \cite[Theorem 3.16]{Ran2021}
gives
\[
[\h_p^c(\M),\h_r^c(\M)]_\eta=\h_1^c(\M).
\]
Setting \(\phi=1-1/r\), the Banach endpoint theorem of
Bekjan--Chen--Perrin--Yin \cite[Theorem 4.4]{BCPY2009}, with the couple
reversed, gives
\[
[\h_1^c(\M),\bmo^c(\M)]_\phi=\h_r^c(\M).
\]
Wolff's inclusion lemma, in the quasi-Banach inclusion form used in
\cite[Corollary 3.19]{Ran2021}, then yields
\[
[\h_p^c(\M),\bmo^c(\M)]_\theta\subset \h_r^c(\M),
\]
after the standard parameter calculation
\(\theta=\phi/(1-\eta+\eta\phi)\).

It remains to prove the reverse inclusion under the stated dependency. By
the lemma, the identity map gives a continuous endpoint embedding
\[
(\h_p^c(\M),\h_\infty^c(\M))
\longrightarrow
(\h_p^c(\M),\bmo^c(\M)).
\]
The complex interpolation functor is monotone for such endpoint embeddings,
so
\[
[\h_p^c(\M),\h_\infty^c(\M)]_\theta
\hookrightarrow
[\h_p^c(\M),\bmo^c(\M)]_\theta .
\]
Combining this with the conditional dependency gives
\[
\h_r^c(\M)
\hookrightarrow [\h_p^c(\M),\h_\infty^c(\M)]_\theta
\hookrightarrow [\h_p^c(\M),\bmo^c(\M)]_\theta .
\]
Together with the first inclusion, this proves the conditional equality.
\end{proof}

\begin{theorem}[conditioned endpoint reduction]
Let \(0<p<1\), \(0<\theta<1\), and \(1/r=(1-\theta)/p\), with \(r>1\). Under
the conditioned \(L_\infty\) endpoint interpolation dependency above,
\[
\h_r^c(\M)\hookrightarrow[\h_p^c(\M),\h_\infty^c(\M)]_\theta.
\]
Consequently,
\[
[\h_p^c(\M),\bmo^c(\M)]_\theta=\h_r^c(\M)
\]
with equivalent quasi-norms.
\end{theorem}

\begin{proof}
Let \(x\in\h_r^c(\M)\). Its martingale difference sequence
\(D_cx=(dx_n)\) belongs to \(L_r^{\rm cond}(\M;\ell_2^c)\), with
\[
\|D_cx\|_{L_r^{\rm cond}}=\|x\|_{\h_r^c}.
\]
By the conditioned endpoint dependency, \(D_cx\) belongs to
\[
[L_p^{\rm cond}(\M;\ell_2^c),
  L_\infty^{\rm cond}(\M;\ell_2^c)]_\theta
\]
with controlled norm. Choose an admissible analytic family \(F\) in the
conditioned couple such that \(F(\theta)=D_cx\). Applying the bounded endpoint
projection \(\Pi\) pointwise gives an admissible analytic family \(\Pi F\) in
\((\h_p^c(\M),\h_\infty^c(\M))\). Its value at \(\theta\) is
\[
(\Pi F)(\theta)=\Pi(D_cx)=x.
\]
Thus
\[
x\in[\h_p^c(\M),\h_\infty^c(\M)]_\theta
\]
with norm bounded by the conditioned endpoint constant and the endpoint norms
of \(\Pi\). This proves the \(h_\infty^c\) reverse inclusion. The preceding
theorem then gives the \(bmo^c\)-endpoint equality.
\end{proof}

\section*{Verification notes}

The proof is symbolic and uses no computation. The first new point beyond the
previous one-sided packet is the endpoint monotonicity step. The second new
point is the projection reduction from Hardy martingales to conditioned
sequences. The \(h_\infty^c\) dependency is strictly weaker than the full open
\(h_\infty^c\) complex identity, since only the reverse inclusion is needed.
The conditioned dependency is a concrete sufficient route to that reverse
inclusion.

\section*{Limitations and failed closing routes}

This packet does not prove the \(h_\infty^c\) complex endpoint dependency.
The 2022 paper \cite{Ran2022} proves K-closedness and real interpolation for
\((\h_p^c,\h_\infty^c)\), then explicitly states that the complex interpolation
version remains open. The 2024 quasi-Banach function-space Wolff theorem of
Egert--Kosmala \cite{EK2024} is also not directly applicable: it is a theorem
for quasi-Banach function spaces with lattice structure, and it explicitly
leaves the corresponding general \(A\)-convex quasi-Banach-space reiteration
problem open.

The natural attempt to prove the conditioned endpoint dependency is to use the
Junge embedding
\[
U:L_s^{\rm cond}(\M;\ell_2^c)\to
L_s(\M\overline\otimes B(\ell_2(\mathbb N^2))).
\]
For finite \(0<s<\infty\), the source paper proves the needed complex
interpolation theorem. At \(s=\infty\), however, the most direct common
retraction of the range of \(U\) would require a uniform endpoint column Stein
estimate for the adapted projection
\[
(\xi_n)\mapsto ((\mathcal E_{n-1}\otimes{\rm id})(\xi_n))
\]
on \(L_\infty(\M;\ell_2^c)\). This is precisely the endpoint estimate that is
absent from noncommutative Stein theory; the source papers only use
complementation in the interior range \(1<s<\infty\). Thus the obvious ambient
projection route does not close the problem without an additional endpoint
hypothesis. For finite filtrations one obtains such a projection with constants
depending on the number of levels, but that does not give the uniform theorem
asked for in the source problem.

The range \(r\leq 1\) is not covered. Even with the \(h_\infty^c\) reverse
inclusion, one would still need an inclusion
\[
[\h_p^c(\M),\bmo^c(\M)]_\theta\hookrightarrow \h_r^c(\M)
\]
in that range, and the current Wolff/Banach-endpoint argument only gives this
for \(r>1\).

\section*{Novelty check}

Local indexes were searched for arXiv:2108.06341, \(h_p^c\), \(bmo^c\),
\(h_\infty^c\), conditioned spaces, and endpoint complex interpolation. Web and
arXiv searches found the source paper \cite{Ran2021}, the later
\(h_\infty^c\) K-closedness paper \cite{Ran2022}, Pisier's classical and
triangular endpoint theorem \cite{Pisier1992}, Kalton's Hardy-type
interpolation-stability theorem \cite{Kalton1994}, Rutsky's real
interpolation/AK-stability work \cite{Rutsky2018}, and the 2024
quasi-Banach function-space reiteration paper \cite{EK2024}. No full solution
of the original \(bmo^c\) complex endpoint identity was found. Pisier and
Kalton provide useful analogies but require analytic Hardy algebra or
K\"othe-lattice structure not available for the conditioned martingale
subspaces.

\begin{thebibliography}{9}

\bibitem{Ran2021}
N.~Randrianantoanina,
\emph{Interpolation between noncommutative martingale Hardy and BMO spaces:
the case \(0<p<1\)},
arXiv:2108.06341.

\bibitem{BCPY2009}
T.~N. Bekjan, Z.~Chen, M.~Perrin, and Z.~Yin,
\emph{Atomic decomposition and interpolation for Hardy spaces of
noncommutative martingales},
arXiv:0906.4437.

\bibitem{Ran2022}
N.~Randrianantoanina,
\emph{P. Jones' interpolation theorem for noncommutative martingale Hardy
spaces},
arXiv:2212.08714.

\bibitem{EK2024}
M.~Egert and B.~W. Kosmala,
\emph{A note on complex interpolation of quasi-Banach function spaces},
arXiv:2412.12769.

\bibitem{Pisier1992}
G.~Pisier,
\emph{Interpolation between \(H^p\) spaces and non-commutative
generalizations I},
arXiv:math/9201229.

\bibitem{Kalton1994}
N.~J. Kalton,
\emph{Complex interpolation of Hardy-type subspaces},
arXiv:math/9402206.

\bibitem{Rutsky2018}
D.~Rutsky,
\emph{Real interpolation of Hardy-type spaces and BMO-regularity},
arXiv:1811.10128.

\bibitem{Wolff}
T.~Wolff,
\emph{A note on interpolation spaces},
Harmonic Analysis (Minneapolis, Minn., 1981), Lecture Notes in Math. 908,
Springer, 1982, 199--204.

\end{thebibliography}

\end{document}
